\documentclass[11pt,a4paper]{article}
\usepackage{amsmath,amssymb}
\usepackage{booktabs}
\usepackage{microtype}
\usepackage{parskip}

\title{Cadmium --- Implementation Notes\\
  \large Assumptions and Context Common to All Experiments}
\author{Damián Vicino}
\date{2026}

\begin{document}
\maketitle

This document records the formalism choices, simulator capabilities, and
conventions that apply to every experiment in this repository.
Experiment-specific specs should reference these notes rather than repeat them.

% ─────────────────────────────────────────────────────────────────────────────
\section*{Foundational Reference}

The tick-counter experiment reproduced in this repository originates from:

\begin{quote}
  Damian Vicino, Olivier Dalle, Gabriel A.\ Wainer.
  \textit{A Data Type for Discretized Time Representation in DEVS.}
  SIMUTOOLS 2014. \texttt{hal-01055555}.
\end{quote}

% ─────────────────────────────────────────────────────────────────────────────
\section*{Formalisms Available in Cadmium}

Cadmium implements two related but distinct DEVS-family formalisms.
Each experiment directory states which formalism it targets.

\subsection*{PDEVS — Parallel DEVS}

PDEVS transitions \emph{all} imminent models simultaneously.  At each time
step the set of imminent models $I \subseteq D$ is computed; every model in
$I$ fires its output function $\lambda$ and then its internal transition
$\delta_{\text{int}}$; models not in $I$ that receive routed output undergo
their external transition $\delta_{\text{ext}}$.  When a model is both
imminent and receives external input, the \textbf{confluence function}
$\delta_{\text{con}}$ resolves the two transitions:
\[
  \delta_{\text{con}}(s, e, x) =
    \delta_{\text{ext}}\!\left(\delta_{\text{int}}(s),\; 0,\; x\right)
  \quad \text{(internal-first, the Cadmium default)}
\]

Multiple simultaneous messages are delivered as a \textbf{message bag}
(\texttt{std::vector<T>} per port), allowing many values per port per step.

\subsection*{Classic DEVS — Sequential DEVS}

Classic (sequential) DEVS serialises event processing: at each time step
\emph{exactly one} imminent model transitions.  When more than one model is
scheduled at the same time, the \textbf{SELECT} function breaks the tie:
\[
  \text{SELECT} : 2^D \setminus \{\emptyset\} \to D,\quad
  \text{SELECT}(I) \in I
\]
The chosen model fires; outputs are delivered to influenced models as
\textbf{message boxes} (at most one value per port per step); then time
advances and the next step begins.

\subsection*{Key Differences}

\begin{center}
\begin{tabular}{@{}lll@{}}
  \toprule
  & PDEVS & Classic DEVS \\
  \midrule
  Imminent set   & all simultaneous models & one model (SELECT) \\
  Tie-breaking   & confluence function      & SELECT function \\
  Message carrier & bag (\texttt{vector})  & box (at most one) \\
  Coordinator    & \texttt{pdevs::coordinator} & bespoke per experiment \\
  \bottomrule
\end{tabular}
\end{center}

% ─────────────────────────────────────────────────────────────────────────────
\section*{Typed Port System}

Both formalisms in Cadmium use a \textbf{typed port system}.  Each port is a
distinct C++ type carrying a specific message type; couplings are expressed as
typed \texttt{(sender-port, receiver-port)} pairs verified at compile time.
Port types are declared as:

\begin{itemize}
  \item \texttt{cadmium::out\_port<T>} — output port carrying values of type \texttt{T}
  \item \texttt{cadmium::in\_port<T>}  — input port carrying values of type \texttt{T}
\end{itemize}

Each model declares its ports via \texttt{using input\_ports} and
\texttt{using output\_ports} type aliases.  The compiler rejects any coupling
where the message types do not match.

% ─────────────────────────────────────────────────────────────────────────────
\section*{TIME Type}

Cadmium is parameterised over a \texttt{TIME} template argument that must
satisfy \texttt{cadmium::concepts::Time}: totally ordered, supporting
arithmetic operations, and providing
\texttt{std::numeric\_limits<TIME>::infinity()} for passive time advance.

Common choices and their trade-offs:

\begin{center}
\begin{tabular}{@{}lll@{}}
  \toprule
  Type & Representation & Risk \\
  \midrule
  \texttt{float}  & IEEE 754 single & $1/10$ not exact; accumulation error \\
  \texttt{double} & IEEE 754 double & same issue, smaller magnitude \\
  rational        & exact fraction  & no accumulation error; heavier compute \\
  \bottomrule
\end{tabular}
\end{center}

Each experiment spec states the TIME types it exercises and the expected
impact on simulation correctness.

\end{document}
